\usepackage[french]{babel}
% fziffggzig

%Choix du thème de la police (sans serif, petites capitales dans les titres...) 
\usepackage{amssymb} 
% Des symboles mathématiques
\usepackage{graphicx}
% Pour insérer des images
\usepackage{multicol}
% Pour gérer facilement un environnement à n colonnes
\usepackage{booktabs}
% Pour des beaux tableaux
\usepackage{esvect} 
% Pour de belles flèches sur les vecteurs
\usepackage{esdiff}
% Pour de beaux symboles de dérivation
\usepackage{siunitx}
% Pour l'utilisation des unités du système international
\sisetup{output-decimal-marker={,},inter-unit-product=\ensuremath{{}\cdot{}}}
% Quelques réglages pour l'affichage des unités

% Pour régler les marges de l'article
\usepackage{geometry}
\geometry{hmargin=2cm,vmargin=1.7cm}
% Pour les couleurs
\usepackage{color}
% Pour les indentations
\usepackage{tcolorbox}
\tcbuselibrary{skins, breakable}
\newtcolorbox{indentation}[1][]{
	blanker, left=1cm, borderline west={0.5mm}{0.6cm}{black!50}, breakable,
	before skip=2ex plus 0.1ex, after skip=2ex plus 0.1ex, #1
}
% Commande pour utiliser une indentation :
% \begin{indentation}
%     <texte>
% \end{indentation}